\addcontentsline{toc}{chapter}{Abstract}\vspace{-1cm}
%Border
\begin{tikzpicture}[remember picture, overlay]
  \draw[line width = 4pt] ($(current page.north west) + (0.75in,-0.75in)$) rectangle ($(current page.south east) + (-0.75in,0.75in)$);
\end{tikzpicture}
Automated Python vulnerability detectors are evaluated almost entirely on
labels inherited from Common Vulnerabilities and Exposures (CVE) fix commits,
where empirical audits report 25\% to 50\% label noise from co-changed files. A
benchmark built on such labels measures a detector's agreement with label noise
as much as its detection ability. The CASTLE benchmark addressed this for C and
C++ with a hand-curated corpus, but no comparable evaluation substrate existed
for Python, which is the dominant language for security tooling, machine
learning pipelines, and modern web applications. A second gap compounded the
first: detection accuracy was reported on unmodified code only, so a reported
score was an upper bound that might not survive a trivial, behavior-preserving
rewrite of the input.

This project constructed a benchmark for Python vulnerability detection covering
seven sink-shaped Common Weakness Enumeration (CWE) classes, with 869 samples
made up of 440 audit-validated real-world positives and 429 synthesized safe
hard negatives. Labels were validated by a three-layer methodology: a pre-ingest
sink-presence filter, a three-round adversarial audit by independent reviewers,
and a diff-based filter that requires the sink line itself to change between the
vulnerable and fixed versions. The benchmark was paired with seven
semantics-preserving adversarial mutators across three families, each targeting
one detector shortcut, and seven detectors were evaluated across four tiers:
pattern-matching static analysis, whole-project dataflow, frontier large
language models (LLMs), and a fine-tuned GraphCodeBERT classifier.

The combined validation pipeline reduced the audited false-positive rate from
52\% to 26\%. Across the seven detectors, no single tool led on every measure,
and the sink-token mutator family separated pattern-matching static analysis,
where Bandit lost 0.187 macro-F1 under attribute-access obfuscation,
from a learned detector that was invariant to every mutator in the suite. All
artifacts, audit reports, rejection manifests, and tool runners were released to
support replication.

\pagebreak